\begin{tabularx}{\linewidth}{l p{1cm} p{1.1cm} p{1cm} p{1.1cm} X}
            \toprule
\rotatebox{60}{\textbf{Dataset}}                 & \rotatebox{60}{\textbf{No. Images}}  & \rotatebox{60}{\textbf{Resolution}} & \rotatebox{60}{\textbf{No. Classes}} & \rotatebox{60}{\textbf{Train/Test Domains}} & \rotatebox{60}{\textbf{Description}}\\
\midrule
\textbf{NICO}~\cite{nico}         & 230000    & Varying   & 60    & 1/5   &  domain-adaptation classification dataset partitioned by a number of contexts. \\
\textbf{Office-Home}~\cite{officehome}    & 15500     & Varying         & 65    & 1/3   &  A domain-adaptation classification dataset partitioned by modality\\
\textbf{Office-31}~\cite{office31}      & 4110      & Varying   & 31    & 1/2   & A domain-adaptation classification dataset partitioned by camera type\\
\textbf{CCT}~\cite{cct}            & 243187    & 1024X747         & 15    & 65/65 & a domain-adaptation/generalization benchmark consisting of images taken from different camera-traps. \\
\textbf{Polyps}~\cite{endocv2020, cvc-clinic, kvasir, etis-larib} & 2720 & arying & 1 & 1/3 & Polyp segmentation datasets collected from different centres. We use Kvasir as the \gls{ind} dataset.\\
\bottomrule
\end{tabularx}